\documentclass[conference]{IEEEtran}
\IEEEoverridecommandlockouts
\usepackage{cite}
\usepackage{amsmath,amssymb}
\usepackage{graphicx}
\usepackage{array}
\usepackage{xcolor}
\usepackage{url}
\usepackage[hidelinks]{hyperref}
\graphicspath{{../docs/figures/}}
% Visible marker for anything still to be checked; must be empty before submission.
\newcommand{\todo}[1]{\textcolor{red}{[TODO: #1]}}

\begin{document}

\title{What Benchmark Scores Hide: Auditing a Flow-Based IoT Intrusion Detector from Labels to Live Blocking}

\author{\IEEEauthorblockN{Muhammad Usama}
\IEEEauthorblockA{Department of Computer Science, COMSATS University Islamabad, Pakistan\\
muhammadusamahoyrr@gmail.com}}

\maketitle

\begin{abstract}
Benchmark scores for network intrusion detectors hide problems that appear only in deployment. We audit a decision tree on 54 NFStream flow features from labels to live blocking on an emulated Open vSwitch gateway, and condense the audit into four checks. \emph{Labels:} 91--98.5\% of the flows in the CIC-IoT2023 web, brute-force and MITM captures are background, i.e.\ benign traffic of other devices recorded alongside the named attack. Trained on file-name labels, the tree flags 71--88\% of the background flows in an attack capture it never saw, against 0.5--5.4\% with evidence-based labels. On IoT-23, with the authors' Zeek labels, it flags 23--47\% of unseen benign devices' flows, against 0--17.5\%. \emph{Splits:} on CIC-IoT2023 a random split hides that brute-force recall falls from 1.00 to 0.05 when the target changes from SSH to RTSP; this case, on three host pairs, makes up most of the random--temporal gap, and the rest is small. On IoT-23 the gap comes from classes missing from one half of the split. \emph{Live:} matching live to offline flows one by one exposed a forwarding pitfall, idle flow meters and an NFStream export-ordering defect that raised a window's median verdict delay to 82.5\,s; later, it showed slow replay splitting flows while macro-F1 stayed unchanged. \emph{Response:} blocking a host pair on its first attack verdict caused 60 false blocks in ten minutes of traffic; requiring 10 verdicts per pair within 60\,s windows caused none but missed the Mirai and web attackers. Code, attacker evidence and result files are public.

\end{abstract}

\begin{IEEEkeywords}
intrusion detection, Internet of Things, evaluation methodology, dataset labelling, flow export, CIC-IoT2023
\end{IEEEkeywords}

\section{Introduction}\label{sec:intro}

Machine-learning network intrusion detectors are mostly evaluated on public benchmarks, where many report macro-F1 close to the maximum. Such scores say little about an IoT gateway, which has a few CPU cores, sees traffic in time order, builds flows itself and must act on its verdicts. The gap between benchmark and operation is long known~\cite{sommer2010outside} and well documented for security machine learning~\cite{arp2022dos,apruzzese2023sok}; a recent study reports a random forest with hold-out F1 0.97 whose attack recall fell to 0.0095 in operational validation on a labelled laboratory capture~\cite{tang2026hybrid}.

Each source of this gap is known on its own: label errors in benchmarks~\cite{engelen2021troubleshooting,liu2022error,lanvin2023errors,flood2024smells}, splits that ignore time~\cite{pendlebury2019tesseract,arp2022dos,andresini2021insomnia}, exporters and timeouts that change the flows~\cite{velan2020impact,pinto2025flowexporter}, and detectors that fall short once traffic arrives as a stream~\cite{pekar2025early,zhao2026gidseval}. We do not claim a new failure mode, and we propose no new detection method. Instead we follow one ordinary detector along the whole path: a decision tree over 54 NFStream~\cite{aouini2022nfstream} flow statistics, trained on ten captures of CIC-IoT2023~\cite{neto2023ciciot2023}, with the label and split steps repeated on ten IoT-23 captures~\cite{garcia2020iot23}. At each of four steps (labels, splits, live gateway, response) we ask what a standard evaluation would hide and which check reveals it. The contributions are:
\begin{itemize}
\item \emph{A four-check audit} (Table~\ref{tab:checklist}), implemented as scripts; Checks~1--2 were run on both datasets, Checks~3--4 on CIC-IoT2023. Checks~1 and~2 apply known pitfalls~\cite{arp2022dos,flood2024smells} per capture and host pair; Checks~3 (live/offline flow fidelity) and~4 (false blocks) go past the offline evaluation.
\item \emph{The cost of file-name labels on unseen traffic} (Section~\ref{sec:labels}): to our knowledge the first per-capture background share for CIC-IoT2023, and a leave-one-capture-out test on two datasets.
\item \emph{A split check that finds a change of target, not drift} (Section~\ref{sec:splits}), and does not find a general gap.
\item \emph{A live-fidelity check} (Section~\ref{sec:live}) that exposed an NFStream export-ordering defect, to our knowledge previously undiagnosed (nfstream\#262~\cite{nfstream262}; fix submitted as PR~\#263~\cite{nfstream263}), and flow splitting in our own runs that macro-F1 did not show.
\item \emph{Response scored in false blocks} (Section~\ref{sec:response}), over the number $k$ of verdicts required per host pair.
\end{itemize}

\begin{table}[t]
\centering
\caption{The four-check audit and what it caught. Checks~1--2 were run on CIC-IoT2023 and IoT-23, Checks~3--4 on CIC-IoT2023.}
\label{tab:checklist}
\footnotesize
\setlength{\tabcolsep}{3pt}
\begin{tabular}{>{\raggedright\arraybackslash}p{1.1cm}>{\raggedright\arraybackslash}p{3.4cm}>{\raggedright\arraybackslash}p{3.5cm}}
\hline
Check & Measure & Caught here \\
\hline
1 Labels & Background share per capture; FPR of a file-label model on captures it never saw & 71--88\% of unseen CIC-IoT2023 background flagged (0.5--5.4\% with evidence labels); released CSVs allow only file labels \\
2 Splits & Per-class temporal recall with host-pair counts; class presence in both halves; cluster-unit sensitivity; boundary-spanning share & Brute-force target change on 3 host pairs (CIC-IoT2023); gap from missing classes, not robust to the cluster unit (IoT-23) \\
3 Live & Flow-by-flow match on every run; packet accounting; replay rate; delay vs.\ CPU budget & Forwarding pitfall, idle meters, export-ordering defect; flows split at slow replay \\
4 Response & False blocks per host-pair window and per hour; trade-off over $k$ per class & 60 false blocks in 10\,min at $k=1$; Mirai and web attackers lost at high $k$ \\
\hline
\end{tabular}
\end{table}

\section{Related Work}\label{sec:related}

\emph{Labels.} Label and flow-construction errors in NIDS benchmarks are well documented for CIC-IDS2017 and CSE-CIC-IDS2018~\cite{engelen2021troubleshooting,liu2022error}, and correcting them changes detection scores~\cite{lanvin2023errors}. Flood et al.\ catalogue design smells of benchmark datasets~\cite{flood2024smells}, and models trained on such benchmarks can rely on shortcuts rather than attack behaviour~\cite{jacobs2022emperor}. Garcia and Valeros report that most of the datasets they analysed assign labels ``in files or directory names'' (abstract)~\cite{garcia2023labeling}. For CIC-IoT2023~\cite{neto2023ciciot2023}, Ramkumar et al.\ observed that the captures mix benign and attack traffic and re-filtered them per device by MAC address, but did not quantify the mixing~\cite{ramkumar2024diagnosing}; Vitorino et al.\ re-extract CIC-IoT2023 flows and assume that each capture contains only its named attack~\cite{vitorino2025revisiting}. Kostas et al.\ note that the attack pcaps of IoT-NID also contain benign packets, and that the derived IoTID20 flows label all of them as attacks~\cite[App.~A.3]{kostas2024iotgem}. Background traffic in attack captures is therefore not new; we add its share per capture and the cost of file-name labels on captures the model never saw.

\emph{Splits.} Random splits of time-ordered data are a known pitfall in security machine learning~\cite{arp2022dos,pendlebury2019tesseract,apruzzese2023sok}, time-aware evaluation has been applied to NIDS~\cite{andresini2021insomnia}, and tree classifiers degrade over time on MAWI flows~\cite{schraven2026mawiflow}. On IoT-23, Lyon et al.\ train on earlier months and find performance degrading as new malware types appear~\cite{lyon2026iot23drift}. Traffic-MoE evaluates a per-class chronological split (first 40\% of each class's time span for training, last 40\% for testing) on CIC-IoT2023 and CICIoMT2024, using byte-level flow tokens~\cite{zhou2026trafficmoe}; BRIDGE uses an early/late split in a pooled benchmark~\cite{bridge2026tchnet}; and Rasheed and Alobaedy find that drift causes no consistent decrease in performance but ``specific behaviour changes, such as total blindness, sensitivity inflation and relatively stable detection''~\cite{rasheed2026drift}. IoTGeM trains on telnet and tests on RTSP brute force, by design, and finds that the variants do not generalise~\cite{kostas2024iotgem}. Our split is within each capture, and we trace the gap per class and host pair.

\emph{Live detection and exporters.} Learned detectors on IoT gateways are common~\cite{realguard2022}, as are Mininet/Open vSwitch~\cite{lantz2010mininet,pfaff2015ovs} prototypes that block flagged hosts. Flow export is surveyed by Hofstede et al.~\cite{hofstede2014flow}. Export timeouts delay flow-based detection~\cite{hofstede2013realtime}, the choice of expiration timeouts changes the flow records and can severely affect detection~\cite{velan2020impact}, and the exporter changes detection results~\cite{pinto2025flowexporter,pinto2025hera} and what a model sees~\cite{vitorino2025revisiting}. Models trained on complete flows lose up to 30\% precision and recall, under certain conditions, when tested on partial flows~\cite{pekar2025early}, and a streaming replay of graph-based NIDS found that none of 18 detector--dataset pairs could alert as events arrive~\cite{zhao2026gidseval}. We add a check that a live gateway classifies the flows of the offline evaluation.

\emph{Response.} Blocking via OpenFlow/OVS drop rules is standard in SDN DDoS mitigation~\cite{bawany2017ddos,yan2016sdn}, and a recent SoK reviews how machine-learning SDN defences are evaluated~\cite{orsdemir2026sok}. Few works we found report false-block counts on replayed traffic.

\section{Detector and Evaluation Setup}\label{sec:setup}

\emph{CIC-IoT2023.} We use ten captures~\cite{neto2023ciciot2023}: Benign, DDoS-SYN, DoS-HTTP, Recon-PortScan, MITM-ArpSpoofing, DictionaryBruteForce, SqlInjection, XSS, CommandInjection and Mirai-udpplain (about 5.17\,GB). Files over 300\,MB are trimmed to their first $10^6$ packets. The captures map to eight classes: Benign, BruteForce, DDoS-SYN, DoS-HTTP, MITM-ARP (man-in-the-middle by ARP spoofing), Mirai, Recon and Web (SQL injection, XSS and command injection).

\emph{IoT-23.} We use ten IoT-23 captures~\cite{garcia2020iot23}: seven malware captures (Hide\&Seek, Muhstik, Hakai, two Torii, Mirai, Trojan) and three captures of real, uninfected IoT devices (Philips Hue bridge, Amazon Echo, Somfy smart door lock; of 7-1 only its Somfy-01 part), which we call benign-device captures. The Okiru captures (all at least 1\,GB) are left out. Labels are the dataset's own Zeek \texttt{conn.log.labeled}: each flow is joined to the Zeek connection with the same unordered 5-tuple and the nearest start within $\pm2$\,s, otherwise to the Zeek connection whose span contains the flow's start. Of 756,306 flows, 99.99\% match; the 38 unmatched flows are dropped, and only 270 (0.04\%) overlap Zeek rows of more than one class. The six classes are Benign, PortScan, C\&C, Attack (only in capture 3-1), DDoS and FileDownload.

\emph{Flows and features.} We build flows with NFStream~\cite{aouini2022nfstream} (version 6.6.0), using its statistical features, an idle timeout of 15\,s and an active timeout of 120\,s, with identical settings for both datasets and on the live gateway. The model uses 54 statistical features (packet counts and sizes, inter-arrival times, durations, TCP flag counts per direction). Identifiers (addresses, ports, protocol, timestamps) are excluded, because they let a model memorise the capture~\cite{arp2022dos}. CIC-IoT2023 yields 857,206 labelled flows, 511,527 of them DDoS-SYN (a consequence of trimming).

\emph{Detector.} The detector is a scikit-learn~\cite{pedregosa2011sklearn} decision tree with balanced class weights. Limiting depth and leaf size on validation raised the validation Benign FPR (from 1.3\% to 5.7\% at depth 16, leaf 10), so the \emph{deployed tree} (depth 32, leaf 1; 1,767 nodes, 222\,KiB) is effectively unconstrained. On the CIC-IoT2023 temporal split it scores macro-F1 0.725 (Benign FPR 3.8\%), close to a random forest (0.737, at about 94$\times$ the size and 58$\times$ the single-flow latency) and histogram gradient boosting (0.714, Benign FPR 7.6\%). NSGA-II feature selection~\cite{deb2002nsga2,blank2020pymoo} kept 14 of 54 features at no measurable cost (0.723) but barely shrank the tree (211\,KiB). All results use the deployed tree, retrained with the same settings for each dataset and label set (except the check of the released CSVs in Section~\ref{sec:labels}).

\emph{Splits.} The \emph{temporal} split puts the first 70\% of each capture's flows, by start time, into training (on CIC-IoT2023, 600,050 training and 257,156 test flows). The \emph{random} split draws 70\% at random, stratified by class. Every design choice (tree limits, feature subset, blocking policy) is made on a validation set, the last 20\% of each capture's training period.

\emph{Metrics.} We report per-class F1 and recall with test support, macro-F1 over all classes, and the Benign false-positive rate (FPR), the share of Benign test flows predicted as any attack. A \emph{host pair} is an unordered pair of addresses within one capture. Flows from one host pair are not independent, so 95\% confidence intervals come from a \emph{host-pair bootstrap}: 1,000 resamples of whole host pairs. A flow-level bootstrap~\cite{efron1994bootstrap} gives far too narrow intervals: [0.702, 0.741] for the temporal macro-F1, against [0.568, 0.822] with host pairs. Two splits are compared by a \emph{paired} bootstrap over the union of both test sets' host pairs: each resample draws host pairs with replacement and applies the same draw to both splits (a pair absent from one test set adds nothing to it).

\section{Step 1: Labels}\label{sec:labels}

\emph{What the benchmark hides.} Each CIC-IoT2023 capture is named after one attack, and all its flows are commonly labelled with that name; but the captures were recorded on a network of active IoT devices, so the attack is only part of the traffic. We call the Benign flows inside an attack capture \emph{background}. Garcia and Valeros predict that models trained with one label per PCAP ``may learn to detect operating system benign traffic as malicious'' (Sec.~3.2)~\cite{garcia2023labeling}; we measure this on traffic the model never saw.

\emph{Evidence labels.} We identified the attackers and victims in every attack capture from the packets themselves (Table~\ref{tab:background}): HTTP payloads against the DVWA (Damn Vulnerable Web Application) test site for the web attacks, destination ports and their timing for brute force, top talkers for the floods, Mirai and the port scan, and ARP replies for MITM-ARP. A flow is labelled as an attack only if (i) one endpoint is a confirmed attacker and the other a victim, on the attack's port where known, or (ii) for ARP spoofing, its first packet was sent from or to the attacker's MAC address and the attacker's own IP is not an endpoint, i.e.\ it is a victim's traffic relayed through the attacker. Every other flow is Benign.

\begin{table}[t]
\centering
\caption{Background share (Benign flows under evidence labels) of each attack capture, and the evidence that confirmed the attacker. Addresses are last octets in 192.168.137.0/24.}
\label{tab:background}
\footnotesize
\setlength{\tabcolsep}{3pt}
\begin{tabular}{lrr>{\raggedright\arraybackslash}p{3.85cm}}
\hline
Capture & Flows & Backgr. & Attacker evidence \\
\hline
XSS          &   4,992 & 98.5\% & \texttt{<script>} from .147 \\
CmdInj       &   5,639 & 98.3\% & DVWA \texttt{exec} page from .147 \\
MITM-ARP     &  12,208 & 97.6\% & MAC \ldots:27:d5 claims .1 and .65 \\
SQLi         &   7,383 & 94.8\% & \texttt{SLEEP(5)}, \texttt{UNION SELECT} \\
Mirai        &     653 & 93.7\% & 7 bots to .191 (UDP) \\
BruteForce   &  11,824 & 90.9\% & .65 to SSH :22, then RTSP :554 \\
Recon        & 205,065 & 30.2\% & .178: 1,001 ports on 58 hosts \\
DoS-HTTP     &  51,275 & 23.0\% & .66 to .71/.141:80 \\
DDoS-SYN     & 511,936 &  0.1\% & 6 sources to .245:4070 \\
\hline
\end{tabular}
\end{table}

\emph{Background share.} Background makes up 91--98.5\% of the flows in the web, brute-force and MITM-ARP captures (Table~\ref{tab:background}); overall, file-name and evidence labels differ for 13.4\% of the 857,206 flows.

\emph{The cost of file-name labels on unseen traffic.} We trained the same tree on file-name labels and on evidence labels and measured the Benign FPR. Within a split this is partly circular: the file-label tree is tested on the later background of captures whose earlier background it learned under the capture's name. We therefore also train on nine captures and test on the tenth (leave one capture out, LOCO; Table~\ref{tab:filelabel}). On a held-out attack capture, the file-label tree flags 71.4--87.8\% of the background flows, against 0.5--5.4\% for the tree trained on evidence labels. (Holding out the single Benign capture leaves the file-label tree without any Benign training flow, so its FPR there is 100\% by construction.) Within the temporal split the same tree flags 31.4\% of the Benign capture's own test flows, against 1.0\%. Its pooled FPR of 71.8\% (38,034 of 52,993 Benign test flows) mixes both kinds of traffic: 88.6\% of these false positives are background inside attack captures, which it flags at 86.1\%. Its macro-F1 against evidence labels is 0.396, against 0.725 for the evidence-label tree; against file-name labels the two appear as 0.537 and 0.409, because correct Benign verdicts in attack captures then count as misses. The random split gives the same pattern (pooled FPR 70.2\% against 1.1\%).

\begin{table}[t]
\centering
\caption{The tree trained on file-name labels (file) or on evidence (CIC-IoT2023) or Zeek (IoT-23) labels. Macro-F1 scored against file-name labels (as a file-label study would report) or against evidence/Zeek labels. Benign FPR on background in attack captures and on the benign capture(s). LOCO: range over held-out captures. CIC: CIC-IoT2023. --: not computed.}
\label{tab:filelabel}
\footnotesize
\setlength{\tabcolsep}{2.6pt}
\begin{tabular}{lllcccc}
\hline
 & & Trained & \multicolumn{2}{c}{Macro-F1 vs.} & \multicolumn{2}{c}{Benign FPR} \\
Data & Test & on & file & evid./Zeek & backgr. & benign capt. \\
\hline
CIC & temporal & file  & 0.537 & 0.396 & 86.1\% & 31.4\% \\
            & temporal & evid. & 0.409 & 0.725 & 4.8\% & 1.0\% \\
            & LOCO     & file  & --    & --    & 71.4--87.8\% & (100\%) \\
            & LOCO     & evid. & --    & --    & 0.5--5.4\% & 1.5\% \\
\hline
IoT-23      & temporal & file  & 0.649 & 0.169 & 99.97\% & 4.7\% \\
            & temporal & Zeek  & 0.157 & 0.649 & 3.4\% & 0.18\% \\
            & LOCO     & file  & --    & --    & 9.4--99.95\% & 23.4--46.7\% \\
            & LOCO     & Zeek  & --    & --    & 0--49.1\% & 0--17.5\% \\
\hline
\end{tabular}
\end{table}

\emph{A label error found by the model.} Our first DDoS-SYN labelling, from the first 300,000 packets, listed five sources; error analysis then showed 65,525 ``Benign'' flows predicted as DDoS-SYN, all from a sixth source (.36) that joins later. We re-checked every flood victim's sources over the whole trimmed file; all results use the corrected labels.

\emph{Robustness test on IoT-23.} We use IoT-23 to test the result against labels that are not ours (the authors' Zeek labels), not as evidence of practice on IoT-23. For file-name labels we map each malware family to its main behaviour, e.g.\ Hide\&Seek to PortScan, Torii to C\&C and Mirai to DDoS; this mapping is our choice. The background share (Zeek label Benign) of the malware captures ranges from 5.5\% to 99.9\%. The benign-device captures are the out-of-capture test (Table~\ref{tab:filelabel}). Within the temporal split the file-label tree flags 4.7\% of their flows, against 0.18\% for the tree trained on Zeek labels; under LOCO it flags 23.4--46.7\% of a held-out benign device's flows, against 0--17.5\%. Generalisation to an unseen capture is hard regardless of labels: the Zeek-label tree also flags 49.1\% and 42.0\% of the background in the held-out malware captures 3-1 and 42-1. File-name labels make it much worse. In-capture, the file-label tree flags 97--100\% of the background in every malware capture, a number that mainly shows that it learned each capture's name.

\emph{The benchmark as distributed.} The per-pcap CSV files of the release (from the mirror \texttt{bencorn/CIC-IoT-2023}) hold 39 packet-window features with no label, address or timestamp; the label is the file name. The merged CSVs add a \texttt{Label} column, but it equals the capture's attack name and the rows are shuffled. Neither form can be relabelled from evidence: the released format forces file-name labels. Even scored against file-name labels on a random split, a decision tree and a random forest on these CSVs reach macro-F1 0.749 and 0.803, but BruteForce recall of only 0.43 and 0.41 and Web recall of 0.37 and 0.35.

\section{Step 2: Splits}\label{sec:splits}

\emph{What a random split hides.} Table~\ref{tab:perclass} gives per-class results on CIC-IoT2023. Macro-F1 is 0.867 under the random split (host-pair CI [0.733, 0.875]) and 0.725 ([0.568, 0.822]) under the temporal split. The paired bootstrap puts the difference (random minus temporal) at +0.142 [0.025, 0.287], but almost all of it is one class: BruteForce contributes 78\% of the summed per-class F1 difference. Without BruteForce the difference is +0.036 [0.007, 0.199], and over the six classes with at least five host pairs in both test sets it is +0.030 [0.004, 0.098]. We do not read percentile intervals at coarser cluster units: with one cluster per capture or per attacker, 76--97\% of the resamples miss a class, and the capture-level interval of the temporal score ([0.360, 0.719]) excludes its own point estimate. Compared capture by capture instead, the random split scores higher in 6 of the 9 captures where the scores differ (two-sided sign test $p=0.51$, permutation test $p=0.18$). The claim for CIC-IoT2023 is therefore narrow: a random split hides a change of target in one class; a general gap is small and not established under coarse clustering.

\begin{table}[t]
\centering
\caption{Per-class results of the deployed tree on CIC-IoT2023. $n$ and pairs: temporal test flows and host pairs. Recall: temporal, with 95\% host-pair CI; none is given for one pair or for recall 1.000 in every pair.}
\label{tab:perclass}
\footnotesize
\setlength{\tabcolsep}{3pt}
\begin{tabular}{lrrcccl}
\hline
 & & & \multicolumn{2}{c}{F1} & \multicolumn{2}{c}{Temporal recall} \\
Class & $n$ & Pairs & random & temp. & & host-pair CI \\
\hline
Benign     &  52,993 & 3,424 & 0.992 & 0.974 & 0.962 & [0.901, 0.990] \\
BruteForce &     605 &     3 & 0.988 & 0.100 & 0.053 & [0.000, 0.842] \\
DDoS-SYN   & 153,402 &     6 & 1.000 & 1.000 & 1.000 & every pair 1.000 \\
DoS-HTTP   &  10,110 &     1 & 0.988 & 0.919 & 0.985 & 1 pair, no interval \\
MITM-ARP   &      25 &    13 & 0.101 & 0.018 & 0.160 & [0.000, 0.375] \\
Mirai      &      10 &     6 & 1.000 & 0.900 & 0.900 & [0.571, 1.000] \\
Recon      &  39,819 &    18 & 1.000 & 1.000 & 1.000 & every pair 1.000 \\
Web        &     192 &     8 & 0.865 & 0.889 & 0.917 & [0.804, 0.941] \\
\hline
Macro-F1   & 257,156 &       & 0.867 & 0.725 & & \\
\hline
\end{tabular}
\end{table}

\emph{A change of attack target.} The brute-force capture shows what the gap is. Training holds 475 SSH flows from the attacker (.65) to .67:22. The test period holds the last 38 SSH flows of that pair and 567 RTSP flows to two cameras, .11 and .140 on port 554, whose pairs never appear in training. The deployed tree detects 32 of the 38 SSH flows and none of the 567 RTSP flows: its recall of 0.053 comes entirely from the SSH tail. A random split puts RTSP flows into training and hides that the model learned one service, not ``brute force''. This is a change of target within one capture session, not drift, and it rests on three host pairs; we treat it as a case study.

\emph{Boundary-spanning connections.} A temporal cut by flow start does not cut connections: a long connection yields several flows (active timeout 120\,s). Of the temporal test flows, 36.1\% (92,766 of 257,156) share a 5-tuple with a training flow of the same capture, i.e.\ continue a connection that started in training; we call these connections boundary-spanning. Removing them lowers macro-F1 from 0.725 to 0.604, mostly because Mirai is then left without test flows; over the seven classes that keep test flows the change is 0.700 to 0.690, concentrated in DoS-HTTP (F1 0.919 to 0.851). All 10 Mirai test flows are such continuations, so the Mirai test score is not evidence of generalisation.

\emph{What flow features cannot see.} MITM-ARP reaches precision 0.01 and recall 0.16 on the temporal split. An ARP-spoofing flow is a victim's ordinary traffic relayed through the attacker, and the MAC address, the only distinguishing field, is not a model input.

\begin{table}[t]
\centering
\caption{IoT-23, deployed tree trained on Zeek labels. 95\% CIs from a flow-level and a host-pair bootstrap. Common: macro-F1 over the four classes present in both halves of both splits. Spanning: share of boundary-spanning test flows.}
\label{tab:iot23split}
\footnotesize
\setlength{\tabcolsep}{2.5pt}
\begin{tabular}{lccccc}
\hline
Split & Macro-F1 & Flow CI & Host-pair CI & Common & Spanning \\
\hline
Temporal & 0.649 & [0.648, 0.650] & [0.613, 0.655] & 0.973 & 4.1\% \\
Random   & 0.863 & [0.764, 0.923] & [0.638, 0.939] & 0.956 & 5.8\% \\
\hline
\end{tabular}
\end{table}

\emph{IoT-23: the gap does not replicate.} On IoT-23 the temporal macro-F1 is again lower than the random one (Table~\ref{tab:iot23split}), but because of missing classes: under the temporal split DDoS has no training flows and FileDownload no test flows, so both score 0 by construction. On the four classes present in both halves the temporal split scores slightly higher (paired random minus temporal $-0.017$ [$-0.039$, $-0.009$]), but host pairs are near-singletons here (91.5\% of the temporal test flows; CIC-IoT2023: 0.4\%), and with clusters by originating host or by capture the paired interval includes zero ([$-0.130$, 0.114] and [$-0.078$, 0.005]). Boundary-spanning flows concentrate in C\&C (36.5\% of its temporal and 68\% of its random test flows); removing them leaves the temporal score unchanged but lowers the random common-class macro-F1 from 0.956 to 0.918, through C\&C (F1 0.885 to 0.733). We therefore do not claim that the CIC-IoT2023 gap replicates. What carries over is Check~2: verify that each class is present in both halves, report per-class temporal recall with host-pair counts and its sensitivity to the cluster unit, and report the boundary-spanning share.

\section{Step 3: Live Gateway}\label{sec:live}

\emph{What an offline score hides.} A gateway classifies flows exported while packets still arrive, not flows from a finished file. Equal scores do not show that both see the same flows, so we check fidelity flow by flow.

\subsection{Testbed}\label{sec:testbed}

The testbed runs Mininet~2.3.0~\cite{lantz2010mininet} with Open vSwitch (OVS)~\cite{pfaff2015ovs} under WSL2 Ubuntu~24.04 on a laptop (i5-8350U, no GPU). In the topology h1--s1--gw--s2--h2, h1 replays traffic with tcpreplay at $\times1.0$ timing, and the gateway gw captures on its s1-facing interface. The gateway runs NFStream~6.6.0 in live mode with the offline settings, a single meter and a 64\,MB capture buffer. Exported flows are classified in micro-batches (256 flows or 50\,ms) by the deployed tree; each flow receives one \emph{verdict}, its predicted class. The gateway can install OVS drop rules for both directions of a host pair (Section~\ref{sec:response}).

\emph{Windows and harness.} For each capture we cut one 60\,s window that starts 1\,s after the last training flow started, so all replayed traffic comes from the test period. Boundary-spanning connections are removed on both the offline and the live side; the offline reference is NFStream run on the same window file. After each window the harness waits until the gateway is idle and clears OVS blocks, per-pair verdict counts and ARP state, since the sessions reuse addresses. The first no-limit run (run~1) and the one-core run used an earlier harness without the quiet wait and the reset. A cgroup quota (\texttt{systemd-run}) sets the gateway's CPU budget: no limit, one core or half a core.

\emph{Achieved replay rates.} tcpreplay could not reach the original rate of the heaviest windows: DDoS-SYN (17,424 packets per second, pps) and Mirai (4,782\,pps) were replayed at 52--77\% of it.

\subsection{Fidelity Check}\label{sec:fidelity}

\emph{Flow-by-flow match.} We match live to offline flows by 5-tuple and order of appearance within the 5-tuple, on every run. In run~1 (no CPU limit, ten windows), 148,756 of 148,779 offline flows match one of the 148,884 live flows; per window, 96.62--100\% of the matched flows have identical packet and byte counts, and 97.63--100\% receive the same verdict. Over all seven runs without enforcement, at least 99.98\% of the offline flows match and at least 99.97\% of the matched flows receive the same verdict. Macro-F1 over run~1's windows (not comparable to the test-set score) is 0.804 both offline and live.

\emph{Packet accounting.} The packet ratio (packets in live over offline flows, per window) lies between 0.9919 and 1.0064, and NFStream reported zero kernel drops, including in the DDoS-SYN and Mirai windows.

\emph{Delay.} Inference takes 0.44\,ms per batch (medians; 209 flows per batch). Delay comes from flow export, as expected for flow-based detection~\cite{hofstede2013realtime}: the median verdict arrives about 15\,s (the idle timeout) after a flow's last packet, and the first verdict for an attacker pair 15--35\,s after its first packet.

\begin{table}[t]
\centering
\caption{Live runs, DDoS-SYN window (offline: 142,724 flows). Replay: achieved rate and its share of the original 17,424\,pps. Same pkts: share of matched flows with identical packet counts. Delay: median verdict after flow end. Macro-F1: live, over the three windows tested at every CPU budget (offline 0.914).}
\label{tab:repeats}
\footnotesize
\setlength{\tabcolsep}{2.5pt}
\begin{tabular}{lrrrccc}
\hline
 & \multicolumn{2}{c}{Replay} & Live & Same & Delay & Macro- \\
Run & pps & \% & flows & pkts & (s) & F1 \\
\hline
No limit 1  & 12,271 & 70.4 & 142,752 & 1.000 & 17.8 & 0.9171 \\
No limit 2  & 11,611 & 66.6 & 160,242 & 0.877 & 30.9 & 0.9155 \\
No limit 3  & 10,936 & 62.8 & 184,483 & 0.708 & 34.0 & 0.9175 \\
One core    & 13,434 & 77.1 & 142,725 & 1.000 & 21.5 & 0.9284 \\
Half core 1 & 12,088 & 69.4 & 142,730 & 1.000 & 75.9 & 0.9175 \\
Half core 2 & 11,831 & 67.9 & 160,475 & 0.876 & 97.0 & 0.9172 \\
Half core 3 & 12,401 & 71.2 & 142,729 & 1.000 & 60.3 & 0.9173 \\
\hline
\end{tabular}
\end{table}

\emph{Repeats: macro-F1 hides flow splitting.} We ran the no-limit and half-core settings three times each and the one-core setting once (Table~\ref{tab:repeats}). Every run captured every packet (zero kernel drops, DDoS-SYN packet ratio 1.00), but flow construction followed the replay rate. At 69\% or more of the original DDoS-SYN rate the live flow count matched offline within 0.02\%; at 63--68\% it was 12--29\% higher, and the share of matched flows with identical packet counts fell to 0.71--0.88, because slower replay stretches inter-packet gaps past the idle timeout and splits connections. Macro-F1 did not move: 0.9155--0.9175 on the three common windows in every run but the one-core one (0.9284), and 0.8038--0.8043 over all ten windows of the no-limit runs. This is Check~3 applied to our own runs: the score hides the loss of flow fidelity, and only the flow-by-flow match shows it. The one-core run's higher score comes from one extra, correctly classified Mirai flow at a lower replay rate.

\emph{CPU budget.} The budget affected verdict delay only in the DDoS-SYN window (Fig.~\ref{fig:delay}), but its load spilled over: with one core, a backlog kept the CPU at about 100\% for 45--60\,s after the flood, into the BruteForce window (mean CPU 40.6\% there, against 0.8\% without a limit), whose median delay still stayed at 15.7\,s. With half a core the gateway needed another 110--163\,s after the gap to clear its backlog, still without drops: CPU pressure caused lag, not loss, but the flood never filled the 64\,MB buffer. These comparisons are confounded: tcpreplay shares the laptop CPU with the gateway, so the replay rate depends on the gateway's load; the one-core run used the earlier harness; and half a core was tested on three windows only. We report them as observations from these runs. Peak memory was 980--1,025\,MB, from NFStream and Python, not from the 222\,KiB tree.

\begin{figure}[t]
\centering
\includegraphics[width=\columnwidth]{../docs/figures/live_delay.png}
\caption{Median verdict delay after a flow's last packet, per window and CPU budget (legend ``1/2 core'': half a core, tested on the three heaviest windows only). Dashed: the 15\,s idle timeout.}
\label{fig:delay}
\end{figure}

\subsection{Failure Modes the Check Exposed}\label{sec:defect}

The fidelity numbers above hold only after three problems were fixed. None changed an offline score, since at end of file every remaining flow is flushed.

\emph{A forwarding pitfall.} Replayed frames carry the source MACs of many capture devices. A MAC-learning switch maps all of them to h1's port and never forwards frames addressed to them, so part of the traffic silently never reaches the detector. Packet accounting reveals this; one OVS rule that forwards everything from h1 to gw fixes it.

\emph{Idle meters.} NFStream checks a meter's timeouts only when that meter receives a packet, so idle flows can wait in a meter that receives none. We use one meter and run a 200\,pps keep-alive ping between h1 and h2 on a management subnet that is never scored or blocked.

\emph{An export-ordering defect.} Unpatched NFStream~6.6.0 delivered verdicts tens of seconds after the idle timeout. \texttt{NFCache}, an LRU \texttt{OrderedDict}, moves a key to the most-recently-used end only in \texttt{\_\_setitem\_\_}, but \texttt{consume()} updates an existing flow in place, so the cache stays in creation order. \texttt{meter\_scan()} expires flows from the head and stops at the first non-idle flow. One long-lived active flow, such as the keep-alive ping, therefore blocks the idle expiry of every flow behind it until its 120\,s active timeout. A one-line fix, \texttt{cache.move\_to\_end(flow\_key)} in the update branch, restores LRU order without changing flow boundaries or features, so offline comparisons cannot see the defect. It cut the per-window median delay after flow end from 82.5\,s to 15.6\,s (SQLi window) and from 49.2\,s to 15.5\,s (XSS). In the unpatched XSS run only 118 of 224 flows (64.5\% of the packets) had been exported when the run stopped; on a continuously running gateway they would be exported late, not lost. Pekar and Jozsa map NFStream to CICFlowMeter features and leave the causes of value differences, among them ``the specific mechanisms employed for updating and terminating flows'', out of scope~\cite{pekar2024integrity}. To our knowledge the cause of this defect was previously undiagnosed. We reported it as nfstream\#262~\cite{nfstream262} and submitted the fix with a regression test as PR~\#263~\cite{nfstream263}.\footnote{The symptom was reported without a cause in nfstream \#120 (2022)~\cite{nfstream120}.}

\section{Step 4: Response}\label{sec:response}

\emph{What a per-flow FPR hides.} A per-flow Benign FPR of 3.8\% looks modest in a table; at a gateway that blocks host pairs, it means blocked innocent hosts. Garcia and Valeros note that a few false positives can have strong consequences once detections drive IP blocking and that earlier alerts help only at a controlled FPR (Sec.~6.3--6.4)~\cite{garcia2023labeling}, and Akbari et al.\ that, because detectors are prone to false positives, simply dropping suspected traffic can have severe consequences~\cite{akbari2020atmos}. We therefore score response in false blocks, counted per \emph{host-pair window}, a host pair within one 60\,s test window. A block is false if the blocked pair carries no attack flow in that window; the pairs seen in a window without an attack flow are the eligible benign host-pair windows.

\emph{ARP watcher.} For the attack that flow features cannot see (Section~\ref{sec:splits}), a small watcher raises a high alert when a MAC other than the known one announces the router IP (.1), the classic arpwatch idea~\cite{arpwatch}. In the MITM-ARP window it raised one high alert 1.7\,s into the window, and none in the other nine windows. This shows that the rule works here, not a false-alarm rate.

\emph{Policies.} With enforcement on, the gateway blocks a host pair (OVS drop rules for both directions) once the pair has collected $k$ attack verdicts; the harness resets these counters after each window. Optionally, it ignores MITM-ARP flow verdicts, which the ARP watcher covers. Requiring $k$ verdicts is our design choice. We meant to select the policy on validation data, as the one with the fewest false blocks among the ten policies ($k\in\{1,2,3,5,10\}$, MITM-ARP verdicts ignored or not) within 0.005 of the best share of attack flows that start after their pair is blocked. The rule was degenerate: all ten policies lie within 0.0012 of the best share, so it reduced to the fewest false blocks and picked the largest $k$ in the grid ($k=10$, MITM-ARP ignored; 1 false block on validation, against 128 for $k=1$). The result is therefore the trade-off over $k$ (Table~\ref{tab:kcurve}); the enforced policy is one point on it.

\begin{table}[t]
\centering
\caption{Trade-off over $k$, simulated per 60\,s window on the verdicts logged in run~1. Ign.: MITM-ARP verdicts ignored. False: false blocks; /h: per hour of capture time. Class columns: attacker pairs blocked, of the number in parentheses (BF BruteForce, DS DDoS-SYN, DH DoS-HTTP, MA MITM-ARP, Mi Mirai, Re Recon); Att.: their sum. Time: median time from a blocked attacker pair's first flow to its block.}
\label{tab:kcurve}
\footnotesize
\setlength{\tabcolsep}{2.4pt}
\begin{tabular}{rcrrcccccccrr}
\hline
 & & & & BF & DS & DH & MA & Mi & Re & Web & Att. & Time \\
$k$ & Ign. & False & /h & (1) & (6) & (1) & (7) & (7) & (1) & (3) & (26) & (s) \\
\hline
1  & no  & 64 & 384 & 1 & 6 & 1 & 0 & 7 & 1 & 3 & 19 & 19.0 \\
1  & yes & 37 & 222 & 1 & 6 & 1 & 0 & 7 & 1 & 3 & 19 & 19.0 \\
2  & yes &  5 &  30 & 1 & 6 & 1 & 0 & 6 & 1 & 3 & 18 & 33.8 \\
3  & yes &  1 &   6 & 1 & 6 & 1 & 0 & 1 & 1 & 3 & 13 & 15.8 \\
5  & yes &  0 &   0 & 1 & 6 & 1 & 0 & 0 & 1 & 1 & 10 & 15.3 \\
10 & yes &  0 &   0 & 1 & 6 & 1 & 0 & 0 & 1 & 0 &  9 & 15.3 \\
20 & yes &  0 &   0 & 0 & 6 & 1 & 0 & 0 & 1 & 0 &  8 & 15.3 \\
\hline
\end{tabular}
\end{table}

\emph{The trade-off.} Table~\ref{tab:kcurve} is simulated on a run without enforcement, so its $k=1$ counts differ slightly from the enforced run; per-hour rates extrapolate from ten minutes of traffic. Without ignoring MITM-ARP verdicts, false blocks fall from 384 per hour at $k=1$ to 84, 12 and 0 at $k=2$, 3 and 5; ignoring them alone cuts the $k=1$ count from 64 to 37. The DDoS-SYN, DoS-HTTP and Recon attackers are blocked at every $k$, so $k$ trades false blocks against the low-rate attackers. Most Mirai bots are lost at $k=3$ (1 of 7 blocked) and all from $k=5$; the web attackers drop to 1 of 3 at $k=5$ and to none from $k=10$; the brute-force attacker is lost at $k=20$. No MITM-ARP pair is blocked at any $k$. These losses hold within 60\,s windows with per-window counters, where each run saw only 4--5 live Mirai flows; with counters kept over a longer horizon, both attackers and benign pairs would reach a given $k$ more often.

\begin{table}[t]
\centering
\caption{Enforced live blocking over the ten test windows (final harness, blocks cleared per window). Ign.: MITM-ARP verdicts ignored. Attackers: attacker pairs blocked. False: false blocks of eligible benign host-pair windows, and per hour of capture time. Dropped: packets.}
\label{tab:blocking}
\footnotesize
\setlength{\tabcolsep}{3pt}
\begin{tabular}{lrrrrr}
\hline
Policy & Blocked & Attackers & False & /h & Dropped \\
\hline
$k=1$ & 79 & 19/26 & 60/1,674 & 360 & 326,268 \\
$k=10$, ign. & 9 & 9/26 & 0/1,675 & 0 & 163,546 \\
\hline
\end{tabular}
\end{table}

\emph{Enforced runs.} Blocking on the first verdict (Table~\ref{tab:blocking}) blocked all 7 Mirai bots (about 100k packets, 35\% of the window, dropped) but also 60 of 1,674 eligible benign host-pair windows (3.6\%), i.e.\ 360 false blocks per hour of capture time, involving 90 distinct IPs, in 9 of the 10 windows. Most fell in the DDoS-SYN (12.5\%) and DoS-HTTP (6.8\%) windows; 21 were triggered by MITM-ARP verdicts, 14 of them outside the MITM-ARP window. Each cuts a benign pair off for the rest of the window. The $k=10$ policy blocked all six DDoS attackers (133k packets dropped), the DoS-HTTP attacker (30k), the brute-force attacker (217) and the port scanner (0; the scan had ended), with no false blocks, and missed the 7 Mirai bots and the 3 web attackers. In the $k=1$ run the DDoS-SYN window was replayed at only 52\% of the original rate; since slower replay splits flows (Section~\ref{sec:fidelity}), its flow and block counts in that window may be inflated.

\emph{Late blocks.} In the run-1 simulation only 11.1\% of the attack flows start after their pair is blocked, at every $k$: a block needs a verdict, and a verdict needs a flow to end. Blocks still cut the later packets of long attack flows: the enforced runs dropped 0.503 ($k=1$) and 0.342 ($k=10$) of the DDoS-SYN window's packets and 0.707 and 0.698 of the DoS-HTTP window's. The validation share (0.995) is thus far too optimistic, as validation ignores the export delay.

\section{Discussion}\label{sec:discussion}

\subsection{Threats to Validity}\label{sec:threats}

\emph{Internal.} On CIC-IoT2023 the same author identified the attackers, wrote the labelling rule and made all design choices. The labels were revised once after error analysis (the sixth DDoS source), towards the model's predictions; we re-checked every flood victim over the whole file and release the evidence. Our CIC-IoT2023 results are scored against our own labels; IoT-23 supports the cost of file-name labels with the authors' labels but does not validate ours, and its file-name labels depend on our family mapping.

\emph{Construct.} Macro-F1 weighs tiny classes (Mirai, MITM-ARP) like large ones and, on IoT-23, missing classes; we report per-class results with host-pair counts. False blocks per hour extrapolate from ten minutes of traffic and depend on the per-window counter reset.

\emph{External.} Labels and splits were audited on two datasets, but the live and response steps only on CIC-IoT2023, with one capture session per attack class, so they measure generalisation within a session, not across networks. We do not test across datasets, where results are known to be asymmetric~\cite{ali2026siamxbert}. Large CIC-IoT2023 captures were trimmed, so a temporal split covers only their start. Live results hold only at the achieved replay rates, on a laptop where tcpreplay shares the CPU with the gateway, and for floods that did not fill the capture buffer. Replay is open-loop, so ``packets dropped'' is counterfactual: drops cause no retransmissions, and attackers do not adapt.

\emph{Statistical conclusion.} The one-core and enforced settings ran once. The SSH-to-RTSP change rests on three host pairs, and the CIC-IoT2023 gap beyond it is small and not established at capture level. LOCO tests one held-out capture at a time, and its ranges span ten captures per dataset. ``No packets were lost'' refers only to these runs.

\subsection{Conclusion}

Each step of the audit exposed a problem that a standard evaluation hid, from labels to blocking, and once in our own live results, where flow splitting left macro-F1 unchanged. The four checks that caught them are the reusable result. Future work includes a per-class criterion for $k$ with persistent counters, longer live runs, a second live dataset, early classification~\cite{pekar2025early} and edge hardware.

\subsection{Artifact Availability}

The repository \url{https://github.com/muhammadusamahoyrr/EdgeFlow-IDS} (MIT licence) holds the code, attacker evidence, configuration, result files and audit scripts. The scripts recompute the offline numbers from the datasets; the live numbers also need the live flow logs, which are not in the repository. All random choices (split, trees, NSGA-II, bootstrap) use seed 0; live runs are timing-dependent. The NFStream fix is included (\nolinkurl{src/patch_nfstream.py}, diff in \nolinkurl{patches/}).
% PLACEHOLDER for the author (do not compile in until minted): add one sentence naming the tagged
% release used for this paper and its Zenodo DOI, e.g. "The exact version used here is release
% <tag>, archived at Zenodo (DOI <to be minted>)."
The CIC-IoT2023 and IoT-23 data are available from their authors~\cite{neto2023ciciot2023,garcia2020iot23}.


\section*{Acknowledgment}
AI assistants (Anthropic Claude) were used to help write the experiment code and to draft and
edit this text. The author directed the work, checked the experiments and their outputs, and takes
full responsibility for the content.

\bibliographystyle{IEEEtran}
\bibliography{refs}

\end{document}
